\begin{resumo}[Abstract]
 \begin{otherlanguage*}{english}
   The growing digitalization of healthcare services and the overload of care systems have driven the search for automated solutions for clinical triage and anamnesis. However, the direct application of Large Language Models (LLMs) in this domain faces critical challenges related to factual hallucination and the lack of structure in the generated data. This work proposes and technically evaluates a software architecture based on the Backend for Frontend (BFF) pattern for assisted anamnesis. The solution adopts a hybrid Artificial Intelligence strategy, orchestrating the generative model Llama-3.1 for empathetic dialogue conduct and the extractive model Biomedical-NER for the identification and categorization of clinical entities, such as symptoms and medications. The methodology involved literature review, requirements definition, incremental prototyping, and the implementation of a Node.js API integrated with a reactive web interface. Preliminary results from the functional prototype indicate the technical feasibility of the approach, showing that the separation of responsibilities between text generation and data extraction reduces inconsistencies and yields structured outputs for future integration with Electronic Health Records (EHR). The work establishes the architectural basis for the next project stage, including broader testing and validation with users.

   \vspace{\onelineskip}
 
   \noindent 
   \textbf{Key-words}: Artificial Intelligence in Healthcare. Assisted Anamnesis. Software Engineering. LLMs. NER.
 \end{otherlanguage*}
\end{resumo}
